\documentclass[11pt,letterpaper]{article}

\usepackage[spanish,mexico,english]{babel}
\usepackage[utf8]{inputenc}
\usepackage[T1]{fontenc}
\usepackage[top=2.4cm,bottom=2.4cm,left=2.5cm,right=2.5cm]{geometry}
\usepackage{amsmath,amssymb}
\usepackage{graphicx}
\usepackage{booktabs}
\usepackage{float}
\usepackage{algorithm}
\usepackage{algpseudocode}
\usepackage{fancyhdr}
\usepackage{xcolor}
\usepackage{parskip}
\usepackage{hyperref}
\hypersetup{colorlinks=true,linkcolor=black,citecolor=black,urlcolor=blue}

\definecolor{uanlBlue}{RGB}{0,56,101}

\pagestyle{fancy}
\fancyhf{}
\renewcommand{\headrulewidth}{0.4pt}
\fancyhead[L]{\small\textcolor{uanlBlue}{MCIE -- Control No Lineal}}
\fancyhead[R]{\small\textcolor{uanlBlue}{Ene-Jun 2026}}
\fancyfoot[C]{\thepage}

\begin{document}

% ---------- PORTADA ----------
\thispagestyle{empty}
\begin{center}
  {\bfseries Universidad Autónoma de Nuevo León}\\
  Facultad de Ingeniería Mecánica y Eléctrica\\[0.6em]
  \textcolor{uanlBlue}{\rule{\linewidth}{1pt}}\\[0.6em]
  Maestría en Ciencias de la Ingeniería Eléctrica\\
  Unidad de Aprendizaje: \textbf{Control No Lineal}\\[1.2em]
  {\LARGE\bfseries Actividad Fundamental 2}\\[1em]
  \textcolor{uanlBlue}{\rule{0.6\linewidth}{0.4pt}}\\[1em]
  \begin{tabular}{rl}
    Estudiante: & Roberto Treviño Cervantes \\
    Matrícula:  & 1915003 \\
    Maestro:    & Dr.\ Alberto Cavazos \\
  \end{tabular}
\end{center}

\vfill
\newpage
\setcounter{page}{1}

\selectlanguage{english}

\section*{Multi-body architecture}

Al-Rawashdeh et al.\ propose modeling the lithography scanner as a combination of three open-loop serial chains: the reticle chain ($j=1$, blue), the optics chain ($j=2$, purple), and the wafer chain ($j=3$, red), the latter being of primary interest for the step-and-scan motion problem studied here.

The base frame serves as the common structural root for all three chains. In the reticle chain, a dual-stage architecture is employed to satisfy the competing requirements for travel range and precision. The long-stroke motor (LS) of the planar moving-coil type consists of four coil actuators arranged in a platform shape which thus facilitates large-range travel (of at least 300\,mm) with moderate micrometer accuracy in 6-DOF. However, being directly mounted to the base frame, it is inherently susceptible to vibrations propagated from the rest of the machine.

To achieve nanometric tracking accuracy, a short-stroke motor (SS) is appended to the long-stroke carrier. Since travel in this stage is limited up to 1\,mm, it is realized with Lorentz actuators which provide contact-free, high-bandwidth positioning. Given that the Lorentz force depends entirely on the current within the actuator and is independent of the stator's motion, it effectively decouples the end-effector from the base frame disturbances through magnetic suspension. To minimize linearity errors and optimize force transfer, the long-stroke stage carrying the stator must maintain the optimal magnetic gap to the magnets of the rest of the short-stroke actuator. This arrangement is replicated for the wafer chain.

The optics chain is attached to the base frame through a metrology frame which is in turn supported by pneumatic vibration isolators. While the optics box includes adaptive optics with active elements, it is reduced in this framework to a single lens element. This allows the model to focus on machine compliance and minimizing structural vibration transfer rather than involving active optical control.

The physical connections between these stages, which constitute the primary transmission path for vibration, are represented by lumped-parameter spring-damper systems, according to the Kelvin-Voigt model. As its proponents state, this low-fidelity modeling choice is sufficient to study the dynamics between the multiple bodies during step-and-scan motion.

\section*{Generalized coordinates}

Two distinct joint types arise: \emph{actuated} stages ($i \in \{2,4\}$), which are driven by servo motors and carry a near-infinite effective stiffness that keeps $f_{i_j}$ locked to the nominal reference; and \emph{flex} stages ($i \in \{3,5,7\}$), which are passive compliant connections whose stiffness $K_{i_j}$ and damping $C_{i_j}$ are physical material properties of the mechanical coupling.

The full 6-DOF model is here restricted to in-plane motions, since rotations about the $x$ and $y$ axis, as well as translation on $z$ are used mainly for focusing. The generalized coordinate used to describe the configuration of the $i$-th stage within the $j$-th chain is
\begin{equation}
    q_{i_j} = \begin{bmatrix} f_{i_j} \\ \theta_{z_{i_j}} \end{bmatrix} \in \mathbb{R}^{3\times1},
\end{equation}
where $f_{i_j} = [f_{x,i_j},\, f_{y,i_j}]^T \in \mathbb{R}^{2}$ collects the in-plane translational displacements of stage $i$ relative to its parent body in chain $j$, and $\theta_{z_{i_j}} \in \mathbb{R}$ is the corresponding rotation about the out-of-plane $z$-axis. The joints between bodies within a chain are modeled using a spring-damper system. The associated generalized mass-inertia matrix for each body is
\begin{equation}
    M_{i_j} = \begin{bmatrix} m_{i_j}\,\mathbf{I}_2 & 0 \\ 0 & I_{z_{i_j}} \end{bmatrix} \in \mathbb{R}^{3\times3},
\end{equation}
and the generalized stiffness and damping matrices $K_{i_j},\, C_{i_j} \in \mathbb{R}^{3\times3}$ are block-diagonal, with a $2\times2$ translational block and a scalar rotational entry.

\section*{Kinematics}

The absolute position $f^{0}_{i_j}$ of body $i$ is expressed as the recursive sum of its relative displacement from its parent body along the chain from the floor frame $\mathcal{F}_{0}$. For bodies with $i \geq 2$ on any chain $j \in \{1,2,3\}$:
\begin{equation}
    f^{0}_{i_j} = f^{0}_{(i-1)_j} + f_{i_j}.
\end{equation}
The angular position recurrence is additive:
\begin{equation}
    \theta^{0}_{z_{i_j}} = \theta^{0}_{z_{(i-1)_j}} + \theta_{z_{i_j}}.
\end{equation}
The translational velocity expression is
\begin{equation}
    \dot{f}^{0}_{i_j} = \dot{f}^{0}_{(i-1)_j} + \dot{f}_{i_j} + \omega_{(i-1)_j} \times f_{i_j},
\end{equation}
where $\omega_{(i-1)_j}$ is the angular velocity of $\mathcal{F}_{(i-1)_j}$ with respect to the floor $\mathcal{F}_{0}$. Taking the second derivative gives:
\begin{equation}
    \ddot{f}^{0}_{i_j} = \ddot{f}^{0}_{(i-1)_j} + \ddot{f}_{i_j}
    + 2\!\left[ \omega_{(i-1)_j} \times \dot{f}_{i_j} \right]
    + \dot{\omega}_{(i-1)_j} \times f_{i_j}
    + \omega_{(i-1)_j} \times \!\left[ \omega_{(i-1)_j} \times f_{i_j} \right].
\end{equation}

We restrict the angular velocity of each body purely about the $z$-axis as $\omega = \dot{\theta}_z\,\hat{z}$, and substitute the cross-products by matrix-vector products with the skew-symmetric matrix:
\begin{equation}
    [\omega_{(i-1)_j}] = \dot{\theta}^{0}_{z_{(i-1)_j}}
    \begin{bmatrix} 0 & -1 \\ 1 & 0 \end{bmatrix},
\end{equation}
such that $\omega_{(i-1)_j} \times f_{i_j} = [\omega_{(i-1)_j}]\,f_{i_j}$, and applying this matrix twice gives the centripetal term $[\omega_{(i-1)_j}]^2 = -(\dot{\theta}^{0}_{z_{(i-1)_j}})^2 \mathbf{I}_2$. Expanding recursively from the floor frame, the translational kinematics become:
\begin{align}
    f^{0}_{i_j}        &= f^{0}_{1} + \sum^{i}_{k=2} f_{k_j}, \\
    \dot{f}^{0}_{i_j}  &= \dot{f}^{0}_{1} + \sum^{i}_{k=2} \dot{f}_{k_j}
                          + \sum^{i}_{k=2} [\omega_{(k-1)_j}]\, f_{k_j}, \\
    \ddot{f}^{0}_{i_j} &= \ddot{f}^{0}_{1} + \sum^{i}_{k=2} \ddot{f}_{k_j}
                          + \sum^{i}_{k=2}\!\left[
                              2[\omega_{(k-1)_j}]\dot{f}_{k_j}
                              + [\dot{\omega}_{(k-1)_j}]f_{k_j}
                              + [\omega_{(k-1)_j}]^2 f_{k_j}
                          \right].
\end{align}
Since $\theta_{z_{i_j}}$ is a scalar, no Coriolis or centripetal cross-terms arise in the angular kinematics:
\begin{equation}
    \theta^{0}_{z_{i_j}} = \theta^{0}_{z_1} + \sum^{i}_{k=2} \theta_{z_{k_j}},
    \qquad
    \ddot{\theta}^{0}_{z_{i_j}} = \ddot{\theta}^{0}_{z_1} + \sum^{i}_{k=2} \ddot{\theta}_{z_{k_j}}.
\end{equation}

\section*{Steiner's theorem and inertia aggregation}

The parallel axis (Steiner's) theorem states that the moment of inertia of a rigid body about any axis parallel to one through its center of mass satisfies
\begin{equation}
    I_z = \tilde{I}_z + m\,\|f\|^2,
\end{equation}
where $\tilde{I}_z$ is the body's own moment of inertia about the $z$-axis through its CoG, $m$ is its mass, and $\|f\|^2 = f_x^2 + f_y^2$. In a serial chain, the inertia of every downstream body must be shifted to the frame of the current body. The aggregated inertia for the $j$-th chain is computed recursively from the end-effector inward.

\begin{algorithm}[H]
\caption{Planar inertia aggregation for the $j$-th chain.}
\begin{algorithmic}[1]
\Require{CoG inertias $\tilde{I}_{z_{i_j}}$, masses $m_{i_j}$, relative positions $f_{k_j}$ for $i = 2,\ldots,N_j$}
\Ensure{Aggregated inertias $^0I_{z_{i_j}}$ for $i = 2,\ldots,N_j$}
\State $^0I_{z_{N_j}} \leftarrow \tilde{I}_{z_{N_j}}$
\If{$2 \leq i < N_j$}
    \For{$k = N_j,\;\ldots,\;i+1$}
        \State $\tilde{m}_{(k-1)_j} \leftarrow \sum_{b=k}^{N_j} m_{b_j}$
        \State $^0I_{z_{(k-1)_j}} \leftarrow {}^0I_{z_{k_j}} + \tilde{I}_{z_{(k-1)_j}} + \tilde{m}_{(k-1)_j}\!\left(f_{x,k_j}^2 + f_{y,k_j}^2\right)$
    \EndFor
\EndIf
\end{algorithmic}
\end{algorithm}

The aggregated inertia of the base frame applies the same shift one more time, from body~2's CoG to the base frame:
\begin{equation}
    {}^0I_{z_1} = \tilde{I}_{z_1}
    + \sum_{j=1}^{3}\!\Bigl({}^0I_{z_{2_j}}
      + \tilde{m}_{1_j}\!\left(f_{x,2_j}^2 + f_{y,2_j}^2\right)\Bigr),
\end{equation}
where $\tilde{m}_{1_j} = \sum_{b=2}^{N_j} m_{b_j}$ is the total mass of chain $j$.

\section*{Base-frame equations of motion}

Starting from Newton's equation $F=ma$, for every body ($i\geq 2$) within chain $j$, it connects to body $i-1$ through a spring above and body $i+1$ through the spring below. Summing the resulting equations over all bodies in a chain, the internal spring forces cancel by Newton's third law; only the boundary term connecting body $i=2$ to the base frame survives. The total translational equation of the base frame (absorbing the balance masses and collecting all chain contributions) is:
\begin{equation}
    m_0\,\ddot{f}^{0}_{1}
    + {}^0\bar{C}_1\,\dot{f}_1
    + {}^0\bar{K}_1\,f_1
    - \sum_{j=1}^{3}\!\left\{
        {}^0\bar{C}_{2_j}\,\dot{f}_{2_j} + {}^0\bar{K}_{2_j}\,f_{2_j}
    \right\} = u^0_{1},
\end{equation}
where $m_0 = m_{1} + \sum_{j=1}^{3}\!\left(m^j_{B} + \sum_{i=2}^{N_j} m_{i_j}\right)$ is the total system mass. By the same cancellation applied to the moment balance about the $z$-axis, the rotational equation of the base frame is:
\begin{equation}
    I_{z_0}\,\ddot{\theta}^{0}_{z_1}
    + {}^0\bar{C}_1\,\dot{\theta}_{z_1}
    + {}^0\bar{K}_1\,\theta_{z_1}
    - \sum_{j=1}^{3}\!\left\{
        {}^0\bar{C}_{2_j}\,\dot{\theta}_{z_{2_j}}
        + {}^0\bar{K}_{2_j}\,\theta_{z_{2_j}}
    \right\} = \tau^0_{1}.
\end{equation}

\section*{Equations of motion for each chain body}

For the first body in the chain ($i=2$), which connects to the base frame above and to body $i=3$ below, the translational and rotational equations are:
\begin{align}
    \bar{m}_{2_j}\ddot{f}^{0}_{2_j}
    + {}^0\bar{C}_{2_j}\,\dot{f}_{2_j}
    + {}^0\bar{K}_{2_j}\,f_{2_j}
    - \left\{ {}^0\bar{C}_{3_j}\,\dot{f}_{3_j} + {}^0\bar{K}_{3_j}\,f_{3_j} \right\} &= u^0_{2_j}, \\
    I_{z_{2_j}}\ddot{\theta}^{0}_{z_{2_j}}
    + {}^0\bar{C}_{2_j}\,\dot{\theta}_{z_{2_j}}
    + {}^0\bar{K}_{2_j}\,\theta_{z_{2_j}}
    - {}^0\bar{C}_{3_j}\,\dot{\theta}_{z_{3_j}}
    - {}^0\bar{K}_{3_j}\,\theta_{z_{3_j}} &= \tau^0_{2_j}.
\end{align}

For bodies $3\leq i< N_j$, the same structure holds:
\begin{align}
    \bar{m}_{i_j}\ddot{f}^{0}_{i_j}
    + {}^0\bar{C}_{i_j}\,\dot{f}_{i_j}
    + {}^0\bar{K}_{i_j}\,f_{i_j}
    - \left\{ {}^0\bar{C}_{(i+1)_j}\,\dot{f}_{(i+1)_j} + {}^0\bar{K}_{(i+1)_j}\,f_{(i+1)_j} \right\} &= u^0_{i_j}, \\
    I_{z_{i_j}}\ddot{\theta}^{0}_{z_{i_j}}
    + {}^0\bar{C}_{i_j}\,\dot{\theta}_{z_{i_j}}
    + {}^0\bar{K}_{i_j}\,\theta_{z_{i_j}}
    - {}^0\bar{C}_{(i+1)_j}\,\dot{\theta}_{z_{(i+1)_j}}
    - {}^0\bar{K}_{(i+1)_j}\,\theta_{z_{(i+1)_j}} &= \tau^0_{i_j}.
\end{align}

For body $i = N_j$ (the end effector), there is no subsequent body in the chain:
\begin{equation}
    \bar{m}_{{N_j}_j}\ddot{f}^{0}_{{N_j}_j}
    + {}^0\bar{C}_{{N_j}_j}\,\dot{f}_{{N_j}_j}
    + {}^0\bar{K}_{{N_j}_j}\,f_{{N_j}_j} = 0,
    \quad
    I_{z_{{N_j}_j}}\ddot{\theta}^{0}_{z_{{N_j}_j}}
    + {}^0\bar{C}_{{N_j}_j}\,\dot{\theta}_{z_{{N_j}_j}}
    + {}^0\bar{K}_{{N_j}_j}\,\theta_{z_{{N_j}_j}} = 0.
\end{equation}

\section*{State-space formulation}

Each body possesses three planar degrees of freedom ($f_x$, $f_y$, $\theta_z$) along with their corresponding time derivatives and contributes exactly six states. The complete machine model comprises 6 bodies in the reticle chain ($N_1 = 7$), 3 bodies in the optics chain ($N_2 = 4$), and 6 bodies in the wafer chain ($N_3 = 7$). Including the common base frame (body 1), the system features a total of $1 + (N_1-1) + (N_2-1) + (N_3-1) = 16$ bodies, and the global state vector encompasses $6 \times 16 = 96$ states. Every positional entry within the state vector represents a \emph{relative} joint displacement; the corresponding absolute inertial position $f^0_{i_j}$ is reconstructed using the recursive summations above.

\begin{table}[H]
\centering
\caption{State-space mapping for the global concatenated system vector.}
\begin{tabular}{cll}
\toprule
Variable & Expression & Meaning \\
\midrule
$x_{6p-5}$  & $f_{x_{i_j}}$            & chain $j$, body $i$, $x$-position \\
$x_{6p-4}$  & $f_{y_{i_j}}$            & chain $j$, body $i$, $y$-position \\
$x_{6p-3}$  & $\theta_{z_{i_j}}$       & chain $j$, body $i$, rotation \\
$x_{6p-2}$  & $\dot{f}_{x_{i_j}}$      & chain $j$, body $i$, $x$-velocity \\
$x_{6p-1}$  & $\dot{f}_{y_{i_j}}$      & chain $j$, body $i$, $y$-velocity \\
$x_{6p}$    & $\dot{\theta}_{z_{i_j}}$ & chain $j$, body $i$, angular velocity \\
\bottomrule
\end{tabular}
\end{table}

Solving for the base-frame accelerations yields
\begin{equation}
    \ddot{f}^{0}_{1}
    = -\frac{{}^0\bar{C}_1}{m_0}\,\dot{f}_1
       -\frac{{}^0\bar{K}_1}{m_0}\,f_1
       +\frac{1}{m_0}\sum_{j=1}^{3}\!\left\{
           {}^0\bar{C}_{2_j}\,\dot{f}_{2_j} + {}^0\bar{K}_{2_j}\,f_{2_j}
       \right\}
       +\frac{u^0_{1}}{m_0},
\end{equation}
\begin{equation}
    \ddot{\theta}^{0}_{z_1}
    = -\frac{{}^0\bar{C}_1}{I_{z_0}}\,\dot{\theta}_{z_1}
       -\frac{{}^0\bar{K}_1}{I_{z_0}}\,\theta_{z_1}
       +\frac{1}{I_{z_0}}\sum_{j=1}^{3}\!\left\{
           {}^0\bar{C}_{2_j}\,\dot{\theta}_{z_{2_j}} + {}^0\bar{K}_{2_j}\,\theta_{z_{2_j}}
       \right\}
       +\frac{\tau^0_{1}}{I_{z_0}},
\end{equation}
which together with the chain-body equations above form the global state-space system $\dot{x} = Ax + Bu$ with $A \in \mathbb{R}^{96\times 96}$.

\end{document}
